\chapter{Introduzione}

Negli ultimi anni i modelli di linguaggio di grandi dimensioni (Large Language Models, LLM) hanno assunto un ruolo centrale nello sviluppo dell’intelligenza artificiale. La loro capacità di apprendere strutture linguistiche complesse e conoscenza generale da enormi collezioni di dati, spesso raccolti automaticamente su scala web, ha reso possibile un salto tecnologico che ha trasformato applicazioni di uso quotidiano e sistemi avanzati di produzione del linguaggio. La generazione fluida di testi, le funzionalità di ragionamento, l’automazione di attività complesse e la capacità di adattarsi a compiti estremamente diversi sono ormai caratteristiche riconosciute di questi modelli. Tuttavia, la stessa ricchezza informativa su cui si fondano costituisce anche una delle principali fonti di vulnerabilità.

Durante l’addestramento, un LLM assimila grandi quantità di informazione dai dati disponibili, sviluppando una rappresentazione interna che riflette non solo regolarità linguistiche generali, ma anche contenuti specifici presenti nel dataset. In questo contesto, è naturale distinguere tra due esigenze contrastanti: da un lato, la necessità che il modello \emph{mantenga} (retention) le competenze utili e la conoscenza generale appresa; dall’altro, la possibilità che esso debba \emph{dimenticare} (forgetting) informazioni particolari che risultano indesiderate o problematiche. Sebbene questi concetti possano apparire intuitivi, renderli operativi nei modelli neurali moderni è tutt’altro che banale.

Gli LLM apprendono infatti da dati eterogenei e, inevitabilmente, interiorizzano contenuti che non dovrebbero essere memorizzati: informazioni sensibili, dati personali, materiale protetto da copyright o conoscenze che, in un secondo momento, devono essere rimosse per ragioni legali, etiche o di sicurezza. Un esempio emblematico è il cosiddetto \emph{diritto all’oblio}, sancito da normative come il GDPR, che impone la cancellazione di dati personali su richiesta dell’interessato \cite{feretzakis2025gdpr}. In assenza di meccanismi espliciti di rimozione, i modelli rischiano di continuare a esporre tali informazioni anche dopo che esse dovrebbero essere state eliminate.

Accanto a questi aspetti, la raccolta non filtrata di dati dal web comporta l’assimilazione di numerosi bias. Tra questi rientrano bias demografici (legati a genere, etnia o nazionalità), bias culturali e geografici, stereotipi sociali, nonché bias informativi dovuti alla sovra-rappresentazione di determinate fonti o punti di vista. Inoltre, i dataset possono contenere disinformazione, contenuti tossici o affermazioni fattualmente scorrette, che il modello tende a riprodurre nei propri output \cite{mehrabi2021survey}. Una volta che tali distorsioni vengono apprese, diventa difficile isolarne l’origine e ancora più complesso eliminarne l’impronta dai parametri della rete.

È in questo quadro che si colloca il problema del \emph{machine unlearning}, un'area di ricerca dedicata alla capacità di far ``dimenticare'' selettivamente a un modello una parte dei dati su cui è stato addestrato, preservando al contempo la sua utilità generale. L’idea è intuitiva: ottenere un modello che si comporti come se determinate informazioni non fossero mai state osservate. La formalizzazione offerta da Ginart et al.\ (2019) \cite{ginart2019makingaiforgetyou} stabilisce che un modello può essere considerato realmente ``dimenticato'' di un insieme di campioni soltanto se il suo comportamento risulta indistinguibile da quello che si otterrebbe addestrandolo sin dall’inizio su un dataset privato di tali elementi. Si tratta di un obiettivo teoricamente chiaro, ma complesso da raggiungere in pratica.

La difficoltà deriva dalla natura distribuita dell’informazione nei modelli neurali di grandi dimensioni. L’influenza di un singolo dato non si riflette su un numero ridotto di pesi: si propaga invece attraverso migliaia di parametri, contribuendo a strutture interne che non sono facilmente scomponibili. Per questo motivo l’unlearning non può limitarsi a una rimozione localizzata, ma richiede interventi più articolati, capaci di ridurre l’impatto dei campioni da dimenticare senza compromettere la retention della conoscenza utile.

La soluzione più diretta consisterebbe nel riaddestrare completamente il modello su un dataset ``depurato''. Tuttavia, l’addestramento di un LLM moderno richiede settimane o mesi di calcolo distribuito e risorse computazionali difficilmente accessibili anche per molte aziende tecnologiche \cite{zhao2023survey}. Riaddestrare da zero non è,  quindi, un'opzione praticabile nella maggior parte dei casi. Da qui nasce l’interesse verso tecniche di unlearning \emph{post-hoc}, cioè metodi che intervengono sui parametri del modello già addestrato per ridurre o annullare l’impatto dei dati indesiderati.

Studi più recenti hanno messo in discussione l’affidabilità di molte tecniche di unlearning, suggerendo che alcune procedure possano dare l’illusione di aver rimosso l’informazione, mentre questa rimane ancora recuperabile tramite interrogazioni mirate. Il lavoro di Liu et al.\ (2024) analizza criticamente le metodologie esistenti e mostra come diversi metodi degradino la qualità complessiva del modello senza ottenere una reale eliminazione dell’informazione \cite{liu2024rethinkingmachineunlearninglarge}.

Un ulteriore ostacolo cruciale è rappresentato dalla valutazione: come si stabilisce se un modello ha effettivamente ``dimenticato'' una certa informazione, senza compromettere la sua capacità di rispondere correttamente ad altri compiti? Non esiste un indicatore universale e molte metriche oggi adottate misurano solo approssimazioni del fenomeno \cite{kim2025evaluation, doshi2024blackbox}. Per affrontare queste difficoltà sono nati diversi benchmark concepiti per analizzare aspetti specifici del forgetting. TOFU \cite{maini2024tofu} utilizza dati sintetici per misurare memorizzazione e sensibilità; MUSE \cite{shi2025muse} valuta la rimozione di conoscenze fattuali mantenendo l’equilibrio tra forgetting e retention; WMDP \cite{li2024wmdpbenchmarkmeasuringreducing} si concentra invece sulla capacità dei modelli di ridurre la disponibilità di informazioni considerate “pericolose” attraverso tentativi controllati di estrazione.

Sebbene questi benchmark abbiano contribuito a rendere più rigoroso il confronto tra tecniche, hanno anche messo in evidenza quanto sia difficile distinguere tra un forgetting reale e un deterioramento artificiale delle prestazioni. Per affrontare questa complessità, la piattaforma \emph{OpenUnlearning} \cite{openunlearning2025} propone un quadro unificato per implementare, valutare e confrontare i metodi esistenti, integrando i diversi benchmark e offrendo strumenti utili per analizzare la robustezza delle metriche stesse.

Nonostante i progressi, il machine unlearning rimane un campo ricco di sfide aperte. La necessità di tecniche che scalino ai modelli futuri, la ricerca di garanzie formali di \emph{certified removal}, la protezione da attacchi di extraction sempre più sofisticati e lo sviluppo di metriche più trasparenti rappresentano solo alcune delle questioni centrali per il prossimo futuro. Si tratta di problemi che non riguardano solo la comunità scientifica, ma che incidono direttamente su privacy, sicurezza e trasparenza dei sistemi intelligenti utilizzati ogni giorno.

La presente tesi si propone di investigare un aspetto ancora poco esplorato del machine unlearning: \textbf{l'effetto della variazione linguistica sui processi di unlearning}. In particolare, ci si chiede se l'introduzione di parafrasi durante l'unlearning possa influenzare l'efficacia dei metodi esistenti nel rimuovere informazioni indesiderate preservando al contempo la conoscenza generale e le competenze del modello. Per rispondere a questa domanda, è stato condotto uno studio sistematico su sei tecniche state-of-the-art (DPO \cite{maini2024tofu}, GradDiff \cite{maini2024tofu}, NPO \cite{zhang_negative_2024}, RMU \cite{li2024wmdpbenchmarkmeasuringreducing}, SimNPO \cite{fan_simplicity_2024} e UNDIAL \cite{dong_undial_2024}), analizzando l'impatto di diversi livelli di parafrasi sul benchmark TOFU \cite{maini2024tofu} applicato al modello Llama-3.2-1B-Instruct. L'obiettivo è comprendere se e in che misura la diversità linguistica costituisca un fattore rilevante per ottenere un forgetting più robusto.